% Karteikarten zur Projektpräsentation "Unfallatlas Deutschland (QUA3CK)"
% Eine Karte pro Seite im Format A5 quer (210 x 148,5 mm).
% Für den Druck 2 Karten pro A4-Seite: karteikarten-a4.tex kompilieren.
\documentclass[9pt]{extarticle}

\usepackage[paperwidth=210mm,paperheight=148.5mm,
            left=9mm,right=9mm,top=7mm,bottom=7mm]{geometry}
\usepackage[T1]{fontenc}
\usepackage[utf8]{inputenc}
\usepackage[ngerman]{babel}
\usepackage[autostyle,german=quotes]{csquotes}
\usepackage{microtype}
\usepackage[scale=0.94]{tgheros}
\renewcommand{\familydefault}{\sfdefault}
\usepackage{xcolor}
\usepackage[skins]{tcolorbox}
\usepackage{enumitem}
\usepackage{amssymb}

\tolerance=9999
\setlength{\emergencystretch}{3em}
\hyphenpenalty=200
\pagestyle{empty}
\setlength{\parindent}{0pt}

\definecolor{Accent}{HTML}{315F7D}
\definecolor{AccentDark}{HTML}{24475F}
\definecolor{AccentPale}{HTML}{EAF0F5}
\definecolor{Demo}{HTML}{B85C00}
\definecolor{DemoPale}{HTML}{FDF2E6}
\definecolor{KSI}{HTML}{C22B38}
\definecolor{Grau}{HTML}{5A6672}

% ---------------------------------------------------------------- Kartenkopf
% \karte{Nummer}{Phasenkürzel}{Titel}
\newcommand{\karte}[3]{%
  \clearpage
  {\color{Accent}\rule{\linewidth}{2.4pt}}\par\vspace{2.5pt}
  {\fontsize{8}{9}\selectfont\bfseries\color{Accent}%
   FOLIE #1\hspace{1.2em}\textbar\hspace{1.2em}#2}\par\vspace{1pt}
  {\fontsize{16}{18}\selectfont\bfseries\color{AccentDark}#3}\par\vspace{3pt}
  {\color{Accent!45}\rule{\linewidth}{0.5pt}}\par\vspace{4pt}%
}

% Karte ohne Foliennummer (Cockpit- und Notfallkarten)
\newcommand{\metakarte}[2]{%
  \clearpage
  {\color{AccentDark}\rule{\linewidth}{2.4pt}}\par\vspace{2.5pt}
  {\fontsize{8}{9}\selectfont\bfseries\color{AccentDark}#1}\par\vspace{1pt}
  {\fontsize{16}{18}\selectfont\bfseries\color{AccentDark}#2}\par\vspace{3pt}
  {\color{Accent!45}\rule{\linewidth}{0.5pt}}\par\vspace{4pt}%
}

\newcommand{\rubrik}[1]{%
  \vspace{3.5pt}{\fontsize{7.5}{8.5}\selectfont\bfseries\color{Grau}%
  \MakeUppercase{#1}}\par\vspace{1.5pt}}

% ------------------------------------------------------------------- Boxen
\newcommand{\kernsatz}[1]{%
  \begin{tcolorbox}[enhanced,sharp corners,boxrule=0pt,
    leftrule=2.2pt,colframe=Accent,colback=AccentPale,
    left=5pt,right=5pt,top=3.5pt,bottom=3.5pt,before skip=1pt,after skip=2pt]
    \color{AccentDark}\itshape #1
  \end{tcolorbox}}

\newcommand{\zeigen}[1]{%
  \begin{tcolorbox}[enhanced,sharp corners,boxrule=0pt,
    leftrule=2.2pt,colframe=Demo,colback=DemoPale,
    left=5pt,right=5pt,top=3.5pt,bottom=3.5pt,before skip=3pt,after skip=2pt]
    {\fontsize{7.5}{8.5}\selectfont\bfseries\color{Demo}JETZT ZEIGEN}\par\vspace{2pt}
    #1
  \end{tcolorbox}}

\newcommand{\merksatz}[1]{%
  \vspace{2pt}{\color{KSI}$\blacktriangleright$}\ \textbf{\color{KSI}#1}\par}

% ------------------------------------------------------------------ Listen
\newlist{sag}{itemize}{2}
\setlist[sag]{leftmargin=1.05em,labelsep=0.4em,label=\textcolor{Accent}{\textbullet},
  itemsep=1.3pt,topsep=1pt,parsep=0pt,partopsep=0pt}
\setlist[sag,2]{leftmargin=0.95em,label=\textcolor{Accent!60}{\raisebox{0.25ex}{\rule{0.45em}{0.5pt}}},topsep=1pt,itemsep=0.8pt}

\newlist{begriff}{itemize}{1}
\setlist[begriff]{leftmargin=1.05em,labelsep=0.4em,label=\textcolor{Grau}{\textbullet},
  itemsep=1.1pt,topsep=1pt,parsep=0pt,partopsep=0pt}

\newcommand{\bg}[2]{\item \textbf{#1:} #2}
\newcommand{\frage}[2]{\item \textbf{\color{AccentDark}#1} #2}

% Zweispaltige Kurzlisten auf engem Raum
\usepackage{multicol}

\begin{document}

% ============================================================ COCKPIT-KARTE
\metakarte{Cockpit}{Ablauf, Zeit und Sprungpunkte}

\rubrik{Zeitbudget, 37 Folien in 30 bis 40 Minuten}
\begin{multicols}{2}
\begin{sag}[itemsep=0.9pt]
  \item Einstieg (1-4): 4 Min
  \item Q-Phase (5-8): 5 Min
  \item U-Phase (9-16): 9 Min
  \item A\textsuperscript{3}-Phase (17-24): 9 Min
  \item C-Phase (25-31): 8 Min
  \item K-Phase (32-34): 4 Min
  \item Fazit, Quellen, Dank (35-37): 2 Min
\end{sag}
\end{multicols}

\rubrik{Kürzen erlaubt, niemals kürzen}
\begin{sag}
  \item Kürzbar bei Zeitnot: Folie 11 (Datenqualität), 24 (Threshold), 31 (Contract).
  \item Niemals kürzen: Folie 16 (Machbarkeitswarnung), 20 (arithmetische Decke),
        21 (Reframing), 27 (Champion bleibt).
\end{sag}

\zeigen{%
\begin{sag}[itemsep=1.1pt]
  \item Nach \textbf{14}: Live-Report, U-Phase, Abschnitt 6 \enquote{Bivariate analysis}, Cram\'ers-V-Heatmap. 30 Sek.
  \item Nach \textbf{22}: App, \textit{Model Comparison}, \enquote{Pareto front: macro-F1 vs. Recall(KSI)}. 45 Sek.
  \item Nach \textbf{27}: App, \textit{Model Comparison}, Aufklapper \enquote{Warum bleibt Random Forest der Champion?}. 30 Sek.
  \item Nach \textbf{32}: App, \textit{Risk Predictor}, \enquote{Try an example} anklicken, Vorhersage zeigen. 60 Sek.
  \item Nach \textbf{33}: App, \textit{Overview}, \enquote{Where accidents are disproportionately severe}, Risikobänder umschalten. 45 Sek.
\end{sag}}

\rubrik{Vor dem Start prüfen}
\begin{sag}
  \item Drei Tabs offen: Repo (github.com/jonasyr/unfallatlas-qua3ck),
        Live-Report (jonasyr.github.io/unfallatlas-qua3ck), App (Streamlit Cloud).
  \item App einmal vorher aufrufen. Kaltstart dauert bis zu 30 Sekunden.
\end{sag}

% ==================================================================== TEIL 1
\karte{1}{Einstieg}{Titel}

\kernsatz{Ich habe die Schwere von Verkehrsunfällen in Deutschland aus öffentlichen
Daten vorhergesagt und methodische Sauberkeit wichtiger genommen als eine hohe Zahl.}

\rubrik{Sagen}
\begin{sag}
  \item Name, Titel, Datengrundlage in einem Satz. Nicht die Folie vorlesen.
  \item Drei Live-Artefakte ankündigen: Repository, Live-Report, Streamlit-App.
  \item Ansagen, dass ich an passenden Stellen kurz hineinspringe.
\end{sag}

\rubrik{Begriffe}
\begin{begriff}
  \bg{QUA\textsuperscript{3}CK}{Prozessmodell für ML-Projekte. Question, Understanding,
      Algorithm/Adapt/Adjust, Conclude \& Compare, Knowledge Transfer.}
  \bg{Unfallatlas}{Offener Datensatz der Statistischen Ämter des Bundes und der Länder.
      Jeder polizeilich erfasste Unfall mit Personenschaden seit 2016, eine Zeile pro Unfall.}
\end{begriff}

% --------------------------------------------------------------------------
\karte{2}{Einstieg}{Agenda}

\kernsatz{Die Präsentation folgt exakt dem QUA\textsuperscript{3}CK-Prozess, weil das
Projekt so entstanden ist.}

\rubrik{Sagen}
\begin{sag}
  \item Die fünf Phasen kurz durchgehen, nicht erklären.
  \item Spannungsbogen setzen: \enquote{Die A\textsuperscript{3}-Phase ist die interessanteste,
        weil dort das ursprüngliche Ziel gescheitert ist und ich es begründet ändern musste.}
\end{sag}

% --------------------------------------------------------------------------
\karte{3}{Einstieg}{Das QUA\textsuperscript{3}CK-Prozessmodell}

\kernsatz{QUA\textsuperscript{3}CK zwingt dazu, die Erfolgskriterien festzulegen, bevor man
die Daten kennt, und macht ein Scheitern damit überhaupt erst erkennbar.}

\rubrik{Sagen}
\begin{sag}
  \item Jeden Buchstaben einmal aussprechen.
  \item Kernpunkt: Das Modell ist \textbf{iterativ}. Spätere Phasen dürfen frühere
        Entscheidungen widerlegen.
  \item Betonen: Genau das ist hier passiert. Das ist kein Fehler im Prozess, das ist der Prozess.
\end{sag}

\rubrik{Begriffe}
\begin{begriff}
  \bg{A\textsuperscript{3}}{Drei Tätigkeiten in einer Phase. \textit{Algorithm}: Welche Modellfamilie?
      \textit{Adapt}: Wie passe ich die Daten an, vor allem an das Klassenungleichgewicht?
      \textit{Adjust}: Wie stelle ich die Hyperparameter ein?}
  \bg{Iterativ}{Kein einmaliger linearer Durchlauf, sondern Rücksprung, wenn eine spätere
      Phase eine frühere Annahme widerlegt.}
\end{begriff}

\rubrik{Wenn gefragt}
\begin{begriff}
  \frage{Warum QUA\textsuperscript{3}CK und nicht CRISP-DM?}{Kursvorgabe, inhaltlich sehr ähnlich.
    QUA\textsuperscript{3}CK trennt die drei A-Tätigkeiten explizit, CRISP-DM fasst sie unter
    \enquote{Modeling}. Gemeinsamer Kern: Die Fachfrage steht vor der Modellierung.}
\end{begriff}

% --------------------------------------------------------------------------
\karte{4}{Einstieg}{Motivation}

\kernsatz{Kommunen entscheiden über Verkehrssicherheit meist auf Basis von Unfallzahlen,
nicht auf Basis von Bedingungsprofilen. Genau diese Lücke adressiert das Projekt.}

\rubrik{Sagen}
\begin{sag}
  \item Einstieg mit der Zahl: rund 270.000 Unfälle mit Personenschaden pro Jahr.
  \item Verteilung nennen, sie prägt die ganze Präsentation: \textbf{1 / 18 / 81 Prozent}.
  \item Nutzen: Ein Modell, das aus Bedingungen auf Schwere schließt, kann priorisieren
        statt nur zählen.
  \item Bewusste Ausschlüsse deutlich sagen: keine individuellen Fahrer-Risikoscores,
        keine Versicherungstarifierung. Vorab festgelegt.
\end{sag}

\rubrik{Begriffe}
\begin{begriff}
  \bg{Vision Zero}{Ursprünglich schwedisches, heute EU-weites Ziel: bis 2050 keine
      Verkehrstoten und Schwerverletzten. Der politische Rahmen solcher Analysen.}
  \bg{BASt}{Bundesanstalt für Straßenwesen. Gibt jährlich \enquote{Unfallentwicklung auf
      deutschen Straßen} heraus.}
  \bg{DVR}{Deutscher Verkehrssicherheitsrat.}
\end{begriff}

% --------------------------------------------------------------------------
\karte{5}{Q}{Die Forschungsfrage}

\kernsatz{Die Frage war von Anfang an gestuft formuliert, mit vorab definiertem Ausweichpfad.
Die spätere Zieleänderung ist damit Teil des Designs, keine nachträgliche Rechtfertigung.}

\rubrik{Sagen}
\begin{sag}
  \item Die Frage einmal langsam vorlesen, sie ist lang.
  \item Entscheidender Punkt: Sie startet mit dem Drei-Klassen-Ziel und bleibt explizit offen
        für eine \enquote{evidenzgetriebene operative Umformulierung}.
  \item Vorhergesagt wird aus Bedingungen, die zum Zeitpunkt der Protokollerstellung vorliegen.
  \item In der A\textsuperscript{3}-Phase auf diese Folie zurückverweisen.
\end{sag}

\merksatz{\enquote{Das ist kein nachträglicher Rettungsanker. Das stand in Phase Q im Notebook,
bevor das erste Modell trainiert wurde.}}

\rubrik{Wenn gefragt}
\begin{begriff}
  \frage{Sagt das Modell vorher, ob ein Unfall passiert?}{Nein. Es setzt voraus, dass ein Unfall
    passiert ist, und schätzt, wie schwer er ausgeht. Unfallvorhersage bräuchte Expositionsdaten,
    also Verkehrsmengen. Die hat dieser Datensatz nicht.}
\end{begriff}

% --------------------------------------------------------------------------
\karte{6}{Q}{Hypothesen}

\kernsatz{Drei überprüfbare Hypothesen: zu Zeit und Wetter, zu Ortskontext und zur
gestuften Machbarkeit.}

\rubrik{Sagen}
\begin{sag}
  \item H1 und H2 sind inhaltlich, H3 ist methodisch.
  \item Bei H2 die Einschränkung mitsprechen: Geschwindigkeit, Infrastrukturqualität und
        Rettungszeit sind plausible Mechanismen, werden aber nicht gemessen. Kein Kausalanspruch.
  \item Ergebnis in einem Halbsatz vorwegnehmen:
  \begin{sag}
    \item H1 nur schwach bestätigt, die Wetter-Assoziationen sind sehr klein.
    \item H2 in der schwachen Form bestätigt, Ortsmerkmale tragen messbar bei.
    \item H3 ist der eigentliche Befund der Arbeit.
  \end{sag}
\end{sag}

\rubrik{Begriffe}
\begin{begriff}
  \bg{Hypothese}{Vorab formulierte, überprüfbare Aussage. Entscheidend ist die
      Falsifizierbarkeit: Sie muss durch die Daten widerlegbar sein.}
  \bg{Baseline}{Bewusst einfaches Vergleichsmodell. Schlägt das komplexe Modell die Baseline
      nicht, hat der Aufwand keinen Nutzen.}
\end{begriff}

% --------------------------------------------------------------------------
\karte{7}{Q}{Zielmetriken und Akzeptanzkriterien}

\kernsatz{Macro-F1 als Primärmetrik plus ein zweites Gate auf dem Recall der seltensten Klasse,
damit kein Modell dadurch gut aussieht, dass es die kritischste Klasse ignoriert.}

\rubrik{Macro-F1 von unten aufbaün}
\begin{sag}
  \item \textbf{Precision} einer Klasse: Von allen als diese Klasse vorhergesagten Fällen,
        wie viele waren es wirklich?
  \item \textbf{Recall} einer Klasse: Von allen Fällen, die wirklich diese Klasse sind,
        wie viele hat das Modell gefunden?
  \item \textbf{F1}: Harmonisches Mittel aus beiden. Harmonisch, weil es bestraft, wenn einer
        der beiden Werte sehr klein ist.
  \item \textbf{Macro-F1}: Einfacher Durchschnitt der F1-Werte über alle Klassen. Jede Klasse
        zählt gleich viel, egal wie selten sie ist.
\end{sag}

\rubrik{Sagen}
\begin{sag}
  \item Schlüsselbeispiel: Ein Modell, das immer \enquote{Leichtverletzt} sagt, hat 81 Prozent
        Accuracy. Klingt gut, ist wertlos. Macro-F1 etwa 0,30, Recall auf die tödliche Klasse null.
  \item Selbst Macro-F1 könnte man theoretisch über die beiden häufigen Klassen erreichen.
  \item Das Recall-Gate auf Klasse 1 verhindert genau das.
  \item Auswahldisziplin ansagen: Modellwahl und Schwellenwert nur auf Validierung 2023.
        Test 2024 wird genau einmal geöffnet. \textbf{Wird auf Folie 27 gebraucht.}
\end{sag}

\rubrik{Wenn gefragt}
\begin{begriff}
  \frage{Warum harmonisches Mittel bei F1?}{Das arithmetische Mittel von Precision 1,0 und
    Recall 0,0 wäre 0,5, obwohl das Modell nutzlos ist. Das harmonische Mittel ergibt hier 0.
    Es bestraft Ungleichgewicht zwischen beiden Größen.}
\end{begriff}

% --------------------------------------------------------------------------
\karte{8}{Q}{Rahmenbedingungen und Ethik}

\kernsatz{Interpretierbarkeit war eine harte Vorgabe, keine Kür. Die Ergebnisse informieren
Infrastrukturentscheidungen, nicht Urteile über Personen.}

\rubrik{Sagen}
\begin{sag}
  \item Zügig durchgehen, das ist eine Absicherungsfolie. Nur zwei Punkte betonen.
  \item \textbf{Interpretierbarkeitsvorgabe}: Sie erklärt, warum kein neuronales Netz im
        Kandidatenfeld ist. Ein Verkehrsplaner muss nachvollziehen können, warum ein Ort
        als Risiko markiert wurde.
  \item \textbf{Ethische Rahmung}: Merkmalswichtigkeiten sind Assoziationen, keine Ursachen.
        Dieser Hinweis steht auch in der App selbst.
\end{sag}

\rubrik{Begriffe}
\begin{begriff}
  \bg{Batch-Inferenz}{Vorhersagen werden gebündelt für viele Fälle berechnet, nicht in
      Echtzeit einzeln. Senkt die Anforderungen deutlich.}
  \bg{DSGVO-Konformität quellseitig}{Der Datensatz enthält von vornherein keine
      personenbezogenen Daten. Keine Namen, kein Alter, kein Geschlecht. Die Frage der
      Anonymisierung stellt sich damit gar nicht.}
\end{begriff}

% --------------------------------------------------------------------------
\karte{9}{U}{Der Datensatz im Steckbrief}

\kernsatz{2,09 Millionen Unfälle über neun Jahre, 21 Spalten, öffentlich lizenziert,
vollständig ohne Personenbezug.}

\rubrik{Sagen}
\begin{sag}
  \item Tabelle nicht vorlesen. Drei Zahlen nennen: \textbf{2.092.401 Zeilen, 21 Spalten,
        2016 bis 2024}.
  \item Merkmalsgruppen in einem Durchgang: Zeit, Unfallcharakter, Umgebung, Beteiligung, Raum.
  \item COVID-Beleg, zeigt dass ich die Daten wirklich angeschaut habe: 2020 bricht ein,
        2024 ist mit 268.519 wieder auf dem Niveau von 2019 mit 268.370.
\end{sag}

\rubrik{Begriffe}
\begin{begriff}
  \bg{UKATGEORIE}{Die Zielvariable. Im Quelldatensatz steht tatsächlich \enquote{UKATGEORIE}
      mit Tippfehler statt \enquote{UKATEGORIE}. Die Schreibweise der Quelle wurde bewusst
      unverändert übernommen.}
  \bg{UART vs. UTYP1}{UART ist die \textit{Unfallart}, also die Kollisionsform (Auffahrunfall,
      Zusammenstoß mit Fußgänger). UTYP1 ist der \textit{Unfalltyp}, also die auslösende
      Konfliktsituation (Abbiegeunfall, Fahrunfall). Beide werden nach dem Unfall vergeben.}
  \bg{Parquet}{Spaltenorientiertes Dateiformat. Speichert Spalte für Spalte statt Zeile für
      Zeile, dadurch sehr schnelle Abfragen auf wenige Spalten und starke Kompression.}
  \bg{Git LFS}{Large File Storage. Git-Erweiterung, die große Binärdateien außerhalb der
      Versionsgeschichte speichert und im Repository nur einen Zeiger ablegt.}
  \bg{DuckDB}{Eingebettete analytische Datenbank, konzeptionell wie SQLite, aber
      spaltenorientiert. Führt SQL direkt auf Parquet aus, ohne Import.}
\end{begriff}

% --------------------------------------------------------------------------
\karte{10}{U}{Zielvariable und Imbalance}

\kernsatz{Ein Klassenverhältnis von 1 zu 18 zu 81 macht Accuracy wertlos und führt ohne
Gegenmaßnahme zum Mehrheitsklassen-Kollaps.}

\rubrik{Sagen}
\begin{sag}
  \item Verteilung nennen, sofort die Konsequenz: Ein Modell, das nichts lernt, bekommt
        81 Prozent Accuracy geschenkt.
  \item \textbf{Stabilität über die Jahre}: Die Klassenanteile schwanken über neun Jahrgänge
        um weniger als einen Prozentpunkt. Das rechtfertigt den chronologischen Split. Gäbe es
        starke Drift, wäre Training auf alten Jahren wertlos.
  \item \textbf{Dunkelziffer} aktiv ansprechen, ehrlichster Punkt der Datenphase: Was wir sehen,
        ist auch eine Verteilung des Meldeverhaltens, nicht nur der Realität. Leichte Unfälle
        werden seltener gemeldet. Aus dem Datensatz heraus nicht korrigierbar.
\end{sag}

\rubrik{Begriffe}
\begin{begriff}
  \bg{Klassenungleichgewicht}{Die Klassen sind sehr unterschiedlich häufig. Das Problem ist
      nicht die Ungleichheit selbst, sondern dass Standardverfahren die Gesamtfehlerrate
      minimieren und deshalb die seltene Klasse ignorieren.}
  \bg{Mehrheitsklassen-Kollaps}{Das Modell hat gelernt, immer die häufigste Klasse
      vorherzusagen, weil das die Fehlerrate minimiert.}
  \bg{Dunkelziffer}{Die unbekannte Zahl nie gemeldeter Unfälle.}
\end{begriff}

% --------------------------------------------------------------------------
\karte{11}{U}{Datenqualität und Plausibilisierung \textnormal{\small(kürzbar)}}

\kernsatz{Fünf Klassen von Prüfungen wurden systematisch durchgeführt, und imputiert wurde
bewusst erst später in der Pipeline, nicht hier.}

\rubrik{Sagen}
\begin{sag}
  \item Bei Zeitnot: nur den letzten Punkt machen.
  \item \textbf{Warum in Phase U nicht imputiert wird}: Eine Imputation über den gesamten
        Datensatz zieht Information aus Validierung und Test ins Training. Der Median einer
        Spalte über alle Daten enthält Information aus dem Testjahr 2024. Das ist Leakage.
  \item Konsequenz: Jede Imputation wandert in die sklearn-Pipeline und wird dort pro Fold
        nur auf Trainingsdaten gefittet.
  \item Ehrlichkeitspunkt: 8 Prozent Geokodierungsausfall. Der Datensatz ist bereits
        vorgefiltert, bevor wir ihn bekommen, und diese Filterung ist nicht zufällig.
\end{sag}

\rubrik{Begriffe}
\begin{begriff}
  \bg{Sentinel-Wert}{Ein spezieller Zahlenwert, der \enquote{kein Messwert} bedeutet, etwa
      minus 999. Wird er nicht erkannt, rechnet das Modell mit einer echten Temperatur von
      minus 999 Grad.}
  \bg{Imputation}{Auffüllen fehlender Werte, zum Beispiel mit dem Median.}
  \bg{Selektionsverzerrung}{Die Stichprobe ist nicht zufällig aus der Grundgesamtheit gezogen.
      Hier: Nur geokodierbare Unfälle sind enthalten, und Geokodierbarkeit hängt vermutlich
      mit dem Ortstyp zusammen.}
\end{begriff}

% --------------------------------------------------------------------------
\karte{12}{U}{Anreicherung 1: DWD-Wetterdaten}

\kernsatz{Jeder Unfall wurde über einen k-d-Baum mit der nächsten Wetterstation verknüpft,
mit 99 Prozent räumlicher Abdeckung, aber erheblichen Lücken bei Wind und Sicht.}

\rubrik{Sagen}
\begin{sag}
  \item k-d-Baum erklären: Naiv bräuchte man für 2,09 Mio. Unfälle die Distanz zu rund
        400 Stationen, also fast eine Milliarde Berechnungen. Ein k-d-Baum teilt den Raum
        rekursiv auf, sodass ganze Teilbäume beim Suchen ausgeschlossen werden. Laufzeit
        unter 30 Sekunden.
  \item \textbf{Join-Detail, das oft übersehen wird}: Der Unfallatlas hat kein Tagesdatum,
        nur Jahr, Monat, Stunde. Also wird über alle Tage eines Monat-Stunde-Buckets gemittelt.
        Unfall am 3. Juli um 14 Uhr bekommt das Mittel aller 14-Uhr-Messungen im Juli.
  \item Das ist Rauschen, aber kein Leakage: Die Zukunft fließt nicht ein.
  \item \textbf{Ehrlichkeitspunkt aktiv nennen}: Windgeschwindigkeit fehlt zu 50,8 Prozent,
        Sichtweite zu 54,4 Prozent. Dokumentiert und an A\textsuperscript{3} als Risiko übergeben.
\end{sag}

\rubrik{Begriffe}
\begin{begriff}
  \bg{k-d-Baum}{Datenstruktur zur schnellen Nächste-Nachbarn-Suche in mehrdimensionalen Räumen.}
  \bg{CDC}{Climate Data Center, das Open-Data-Portal des DWD. \textbf{GeoNutzV}: die Verordnung,
      unter der DWD-Daten frei nutzbar sind.}
  \bg{Proxy-Variable}{Ein Merkmal, das für etwas anderes steht, das man nicht direkt messen kann.
      \texttt{dwd\_station\_dist\_km} ist ein Proxy für ländliche Lage. Geprüft: Die Todesrate
      fernab von Stationen liegt über 20 Prozent höher. Deshalb bewusst behalten, obwohl es
      eigentlich ein Datenqualitäts-Artefakt ist.}
\end{begriff}

% --------------------------------------------------------------------------
\karte{13}{U}{Anreicherung 2: OSM-Straßenkontext über H3}

\kernsatz{Statt für 2,09 Millionen Punkte einzeln die nächste Straße zu suchen, wurde das
Straßennetz einmal auf ein Hexagon-Gitter aggregiert.}

\rubrik{Sagen}
\begin{sag}
  \item Erst das Problem: Ein Nächste-Straße-Lookup für 2,09 Mio. Punkte gegen das gesamte
        deutsche Straßennetz ist auf einem Laptop nicht praktikabel.
  \item H3 erklären: Uber-System, zerlegt die Erdoberfläche hierarchisch in Sechsecke.
        Auf Stufe 8 ist eine Zelle etwa 0,7 Quadratkilometer groß.
  \item \textbf{Warum Sechsecke und nicht Quadrate?} Bei Sechsecken haben alle sechs Nachbarn
        denselben Mittelpunktsabstand. Bei Quadraten sind die diagonalen Nachbarn weiter weg
        als die orthogonalen. Für Nachbarschaftsanalysen sauberer.
  \item Der Trick: Straßennetz einmal pro Bundesland laden und auf Zellen aggregieren. Danach
        ist der Join pro Unfall nur noch ein Nachschlagen der Zellen-ID, praktisch kostenlos.
  \item \textbf{Dokumentierte Näherung offen ansprechen}: OSM bildet das heutige Netz ab, die
        Unfälle reichen bis 2016 zurück. Ein Tempolimit kann sich geändert haben. Akzeptierte
        Approximation, kein Leakage, weil OSM-Daten nicht vom Unfallausgang abhängen.
\end{sag}

\rubrik{Begriffe}
\begin{begriff}
  \bg{H3}{Hexagonales hierarchisches Geo-Indexsystem von Uber.}
  \bg{OSM}{OpenStreetMap, offenes von Freiwilligen gepflegtes Kartenprojekt. Lizenz ODbL.}
  \bg{Expositionsproxy}{\texttt{osm\_road\_density} steht indirekt für Verkehrsaufkommen.
      Wo mehr Straße ist, fährt in der Regel mehr Verkehr.}
\end{begriff}

% --------------------------------------------------------------------------
\karte{14}{U}{Zentrale EDA-Befunde}

\kernsatz{Kein einziges Merkmal erreicht auch nur Cram\'ers V von 0,15 zur Zielvariable.
Das schwache Signal ist der zentrale Befund der Datenphase.}

\rubrik{Sagen}
\begin{sag}
  \item Diese Folie ist die Vorbereitung auf das spätere Scheitern. Sorgfältig machen.
  \item Cram\'ers V erklären: Maß für den Zusammenhang zweier kategorialer Variablen,
        abgeleitet aus dem Chi-Quadrat-Test, normiert auf 0 bis 1.
  \item Einordnung liefern: unter 0,1 vernachlässigbar, 0,1 bis 0,3 schwach, 0,3 bis 0,5 mittel,
        darüber stark. \textbf{0,13 ist also am unteren Rand von schwach, und das ist der beste
        Wert im ganzen Datensatz.}
  \item Anschaulichster Teil, die zeitlichen Muster: Die meisten Unfälle passieren im
        Nachmittagsverkehr, die schwersten nachts. Wer zählt, sieht den Nachmittag.
        Wer Schwere betrachtet, sieht die Nacht.
  \item Negativbefund bewusst nennen: Monatlicher Niederschlag und monatliche Todesrate bewegen
        sich nicht gemeinsam. Ein berichtetes Nullergebnis ist ein Qualitätsmerkmal.
\end{sag}

\zeigen{Live-Report $\rightarrow$ \textit{U Phase} $\rightarrow$ Abschnitt \textbf{6 Bivariate
analysis: feature $\times$ target}. Nur die \textbf{Cram\'ers-V-Heatmap} zeigen, auf die
dunkelste Zelle (UART) deuten und die 0,13 vorlesen. Nicht scrollen. 30 Sekunden.}

\rubrik{Begriffe und wenn gefragt}
\begin{begriff}
  \bg{EDA}{Exploratory Data Analysis. Das systematische Anschaün der Daten vor der Modellierung.}
  \frage{Warum nicht Korrelation?}{Pearson setzt metrische Variablen und einen linearen
    Zusammenhang voraus. Unfallart ist nominal, ohne Ordnung. Da ist Korrelation nicht definiert.
    Cram\'ers V ist das passende Gegenstück für kategoriale Variablen.}
\end{begriff}

% --------------------------------------------------------------------------
\karte{15}{U}{Leakage-Prüfung und chronologischer Split}

\kernsatz{Leakage wurde aktiv gemessen, nicht angenommen. Der Split ist chronologisch, weil ein
Zufallssplit bei Zeitreihendaten die Realität verfälschen würde.}

\rubrik{Sagen}
\begin{sag}
  \item Leakage zürst definieren, danach wird am ehesten gefragt: Informationen, die zum
        Vorhersagezeitpunkt nicht verfügbar wären, gelangen ins Training. Folge sind
        Leistungswerte, die im echten Einsatz nicht halten.
  \item Konkrete Sorge hier: UART und UTYP1 werden nach dem Unfall vergeben. Steckt in diesen
        Kategorien die Schwere schon implizit? Dann wäre es ein Zirkelschluss.
  \item Die Sonde, bedingte Entropiereduktion: Entropie misst Unsicherheit. H(Y) ist die
        Unsicherheit über die Schwere ohne Wissen, H(Y\,\textbar\,X) die verbleibende
        Unsicherheit, wenn Merkmal X bekannt ist. Die Reduktion sagt, wie viel Prozent der
        Unsicherheit das Merkmal wegnimmt.
  \item Ergebnis: UART nimmt 3,4 Prozent weg, UTYP1 2,2 Prozent. Weit unter der
        50-Prozent-Auslöseschwelle. Kein Merkmal musste raus.
  \item Split begründen: Bei einem Zufallssplit könnte das Modell auf 2024er-Daten trainieren
        und auf 2019er-Daten getestet werden. Das entspricht nicht dem Einsatz, wo man immer
        in die Zukunft vorhersagt.
  \item Verifikation erwähnen: Keine OBJECTID kommt in mehr als einem Split vor. Geprüft,
        nicht angenommen.
\end{sag}

% --------------------------------------------------------------------------
\karte{16}{U}{Die Machbarkeitswarnung \textnormal{\small(nicht kürzen)}}

\kernsatz{Die Datenphase endete mit einer dokumentierten Warnung, dass das ursprüngliche Ziel
möglicherweise nicht erreichbar ist. Und zwar bevor das erste Modell trainiert wurde.}

\rubrik{Sagen, Wendepunkt der Präsentation, Zeit nehmen}
\begin{sag}
  \item Erstens: Die tödliche Klasse ist etwa 1 Prozent der Daten.
  \item Zweitens: Die stärkste Assoziation ist Cram\'ers V von 0,13, also schwach.
  \item Drittens, der eigentliche Grund, hier langsam werden: Die physikalisch entscheidenden
        Determinanten fehlen komplett. Aufzählen: Aufprallgeschwindigkeit, Alter der Insassen,
        Gurtnutzung, Fahrzeugmasse, Anstoßgeometrie.
\end{sag}

\merksatz{\enquote{Ob ein Unfall tödlich endet, entscheidet sich vor allem an der kinetischen
Energie und daran, wie verletzlich der Mensch im Fahrzeug ist. Genau diese Größen enthält der
öffentliche Datensatz nicht. Was er enthält, ist der Kontext drumherum.}}

\merksatz{\enquote{Diese Warnung wurde dokumentiert, bevor das erste Modell trainiert wurde. Das
ist der Unterschied zwischen einer belegten Vorhersage und einer nachträglichen Ausrede.}}

% --------------------------------------------------------------------------
\karte{17}{A\textsuperscript{3}}{Aufbau der Modellierung}

\kernsatz{Alles läuft in einer sklearn-Pipeline, damit Vorverarbeitung niemals über die
Split-Grenzen leckt.}

\rubrik{Sagen}
\begin{sag}
  \item \textbf{Pipeline}, der methodisch wichtigste Punkt: Verkettet Vorverarbeitung und Modell
        zu einem Objekt mit gemeinsamem \texttt{fit} und \texttt{predict}. Fittet die
        Kreuzvalidierung das Objekt pro Fold neu, werden auch Skalierung, Encoding und Imputation
        nur auf dem Trainingsteil dieses Folds gefittet.
  \item \textbf{Zyklische Kodierung}, lohnt 20 Sekunden: Stunde 23 und Stunde 0 liegen eine Stunde
        auseinander, als Zahlen aber 23. Kodiert man die Stunde als Sinus-Kosinus-Paar auf einem
        Kreis, ist der Abstand wieder korrekt.
  \item \textbf{Target-Encoding}: Statt 87 Kreis-Spalten per One-Hot wird jeder Kreis durch den
        mittleren Zielwert seiner Trainingszeilen ersetzt, mit Glättung gegen Überanpassung
        bei seltenen Kreisen.
  \item \textbf{GroupKFold}: Gruppiert nach Unfalljahr, damit kein Modell im selben Jahr trainiert
        und validiert.
  \item \textbf{Audit-Modus}, starker Punkt: Das Notebook bricht ab, wenn eine persistierte Metrik
        nicht auf 1e-9 genau reproduziert wird. Die Zahlen in dieser Präsentation sind nicht
        abgetippt, sie werden bei jedem Lauf gegengeprüft.
\end{sag}

\rubrik{Begriffe}
\begin{begriff}
  \bg{One-Hot-Encoding}{Eine Kategorie mit k Ausprägungen wird zu k Null-Eins-Spalten.}
  \bg{Kardinalität}{Anzahl verschiedener Ausprägungen einer Kategorie.}
  \bg{Seed}{Startwert des Zufallszahlengenerators. Fixiert man ihn, sind Zufallsentscheidungen
      reproduzierbar.}
\end{begriff}

% --------------------------------------------------------------------------
\karte{18}{A\textsuperscript{3}}{19 Konfigurationen gegen das Drei-Klassen-Ziel}

\kernsatz{Keine der 19 getesteten Konfigurationen erreicht den zulässigen Bereich. Die beste
Macro-F1 liegt bei 0,424 gegenüber der geforderten 0,55.}

\rubrik{Sagen}
\begin{sag}
  \item Tabelle nicht vorlesen, stattdessen das Muster zeigen. Es gibt zwei Gruppen:
  \begin{sag}
    \item Brauchbare Macro-F1, miserabler Recall auf die tödliche Klasse. Oben Random Forest
          balanced mit 0,424 und 0,212.
    \item Brauchbarer Recall, schlechte Macro-F1. XGBoost und LightGBM balanced mit rund 0,60
          Recall, aber nur 0,37 Macro-F1.
  \end{sag}
  \item Einwand vorwegnehmen: Auch eine nachträgliche Verschiebung der Entscheidungsgrenze über
        einen zweidimensionalen Offset-Sweep erreicht das Gate nicht. Es wurde also nicht nur
        \enquote{nicht genug getunt}.
\end{sag}

\merksatz{\enquote{Das Gate verlangt beides gleichzeitig. Und genau die Kombination existiert
nicht. Es gibt keinen Punkt im zulässigen Quadranten.}}

\rubrik{Begriffe}
\begin{begriff}
  \bg{balanced (Klassengewichtung)}{Fehler auf seltenen Klassen werden im Trainingsverlust
      stärker gewichtet, typischerweise umgekehrt proportional zur Klassenhäufigkeit.}
  \bg{Threshold Moving}{Nachträgliches Verschieben der Entscheidungsschwelle, ohne das Modell
      neu zu trainieren.}
  \bg{Offset-Sweep}{Systematisches Durchprobieren additiver Verschiebungen auf den Logits der
      Minderheitsklassen, um zu prüfen, ob irgendein Arbeitspunkt das Gate erreicht.}
\end{begriff}

% ==================================================================== TEIL 2
\karte{19}{A\textsuperscript{3}}{Warum Oversampling das Problem verschlimmert}

\kernsatz{SMOTE und ADASYN lassen den Recall auf die tödliche Klasse auf ein bis drei Prozent
kollabieren. Nur Klassengewichtung wirkt überhaupt.}

\rubrik{Sagen}
\begin{sag}
  \item \textbf{SMOTE} erklären: Synthetic Minority Over-sampling Technique. Statt seltene
        Beispiele zu duplizieren, sucht SMOTE zu einem Minderheitsbeispiel seine nächsten
        Nachbarn derselben Klasse und erzeugt neue Punkte auf den Verbindungslinien dazwischen.
  \item \textbf{ADASYN}: adaptive Variante, erzeugt mehr Beispiele dort, wo die Klassifikation
        schwer ist.
  \item \textbf{Frank-Hall-Zerlegung}: Nutzt aus, dass die Klassen geordnet sind. Statt drei
        Klassen direkt zu lernen, lernt man zwei binäre Fragen: \enquote{schwerer als 1?} und
        \enquote{schwerer als 2?}. Elegant, aber hier mit Recall 0,003 die schlechteste Strategie.
  \item Schluss: Nur Klassengewichtung hebt den Recall über 0,50. Ein Befund, der der Intuition
        widerspricht, denn SMOTE gilt als Standardwerkzeug.
  \item Wichtig: SMOTE darf nur auf Trainingsdaten angewendet werden, niemals auf Validierung
        oder Test. Sonst Leakage.
\end{sag}

\merksatz{\enquote{SMOTE interpoliert zwischen echten Beispielen. Das funktioniert, wenn die
Klassen im Merkmalsraum getrennt liegen. Hier überlappen sie fast vollständig, weil das Signal
so schwach ist. Ein interpolierter Punkt zwischen zwei tödlichen Unfällen landet mitten in einer
Wolke aus Leichtverletzten. Man erzeugt kein Signal, man erzeugt Rauschen.}}

% --------------------------------------------------------------------------
\karte{20}{A\textsuperscript{3}}{Die arithmetische Decke \textnormal{\small(Schlüsselfolie)}}

\kernsatz{Unabhängig vom Modell lässt sich ausrechnen, dass das Gate eine etwa 90-fache
Odds-Anhebung erfordert, die ein Signal von Cram\'ers V 0,13 nicht liefern kann.}

\rubrik{Die vier Schritte langsam durchgehen}
\begin{sag}
  \item \textbf{Schritt 1}: F1 auf der Mehrheitsklasse liegt bei etwa 0,72. Macro-F1 ist der
        Durchschnitt über drei Klassen. Damit er 0,55 erreicht, müssen die anderen beiden
        zusammen im Mittel etwa 0,46 beitragen.
  \item \textbf{Schritt 2}: Klasse 1 macht 0,94 Prozent aus. Für F1 von 0,46 bei Recall
        mindestens 0,50 braucht man Precision von etwa 0,42. Also: Von allen als tödlich
        vorhergesagten Unfällen müssten 42 Prozent wirklich tödlich sein.
  \item \textbf{Schritt 3}, Odds-Rechnung: Odds sind Wahrscheinlichkeit geteilt durch
        Gegenwahrscheinlichkeit. Basisrate 0,94 Prozent ergibt Odds von etwa 0,0095.
        Precision 0,42 ergibt Odds von etwa 0,72. Faktor dazwischen: rund 90.
  \item \textbf{Schritt 4}: Die stärkste Assoziation im Datensatz liegt bei Cram\'ers V von 0,13.
        Ein so schwaches Signal kann diese Konzentration nicht erzeugen.
  \item Abschluss: Empirische Front und arithmetische Anforderung zeigen unabhängig voneinander
        auf dasselbe Ergebnis.
\end{sag}

\merksatz{\enquote{Das Modell müsste eine Teilgruppe finden, in der ein Unfall rund 90-mal
wahrscheinlicher tödlich endet als im Durchschnitt. Nicht 90 Prozent mehr. 90-mal.}}

\rubrik{Wenn gefragt}
\begin{begriff}
  \frage{Ist die 90 exakt?}{Nein, eine Größenordnungsrechnung. Je nach Rundung landet man
    zwischen 76 und 90. Der Punkt ist nicht die zweite Nachkommastelle, sondern die
    Größenordnung: Es braucht eine Konzentration um fast zwei Zehnerpotenzen.}
\end{begriff}

% --------------------------------------------------------------------------
\karte{21}{A\textsuperscript{3}}{Das Reframing auf binäres KSI \textnormal{\small(Schlüsselfolie)}}

\kernsatz{Die Zielvariable wurde auf KSI gegen Leichtverletzt umgestellt. Vorab definiert,
sicherheitsrelevant und mit unverändertem Split.}

\rubrik{Sagen}
\begin{sag}
  \item Entscheidung: Das Drei-Klassen-Ziel bleibt als dokumentierter Hintergrund, ist aber nicht
        länger primär. Begründung in einem Satz: Weiteres Tuning würde dieselbe Suche in
        schwachem Signal wiederholen, ohne die verfügbare Information zu verändern.
  \item \textbf{KSI}: Killed or Seriously Injured, getötet oder schwerverletzt. Etablierter
        Begriff in der Verkehrssicherheitsarbeit, weil er die planerisch relevante Grenze zieht.
  \item Die drei Rechtfertigungen sicher können:
  \begin{sag}
    \item \textbf{Vorab definiert}. Verweis zurück auf Folie 5 und H3. Der Fallback stand im
          Q-Phase-Notebook, bevor modelliert wurde.
    \item \textbf{Sicherheitsrelevant}. Der Unterschied zwischen schwer und leicht verletzt ist
          der Unterschied, an dem sich Maßnahmen entscheiden.
    \item \textbf{Statistisch tragfähig}. Aus 1 Prozent positiver Klasse werden etwa 19 Prozent.
  \end{sag}
  \item Methodische Sauberkeit zeigen: Der zeitliche Split und die Leakage-Grenze bleiben
        unverändert. Es wird nichts nachträglich angepasst außer der Zielvariable selbst.
\end{sag}

\rubrik{Wenn gefragt}
\begin{begriff}
  \frage{Ist das nicht einfach das Problem leichter machen?}{Ja, die Aufgabe wird leichter, und
    das sage ich offen. Der Unterschied zu unsauberem Vorgehen ist dreifach: Der Pfad war vorab
    definiert. Die ursprüngliche Zielverfehlung wird vollständig berichtet, nicht gelöscht.
    Die neue Aufgabe bleibt inhaltlich relevant.}
  \frage{Warum nicht Getötet gegen Rest?}{Das hätte genau das 1-Prozent-Seltenheitsproblem
    behalten. KSI ist die in der Verkehrssicherheitsliteratur übliche Zusammenfassung und löst
    das Häufigkeitsproblem, ohne die planerische Relevanz zu verlieren.}
\end{begriff}

% --------------------------------------------------------------------------
\karte{22}{A\textsuperscript{3}}{Zehn Kandidaten im binären Setting}

\kernsatz{Zehn Modellfamilien von der Zufallsbaseline bis zu drei Boosting-Verfahren, alle auf
identischer Merkmalsbasis und nur auf Validierung verglichen.}

\rubrik{Sagen}
\begin{sag}
  \item Baselines zürst, sie zeigen dass das Modell wirklich etwas lernt: Zufallsraten 0,439,
        Mehrheitsklasse 0,453. Das ist die Latte, alles darüber ist echter Lerngewinn.
  \item \textbf{Random Forest}: Viele Entscheidungsbäume, jeder auf einer Zufallsstichprobe der
        Daten und einer Zufallsauswahl der Merkmale, am Ende wird gemittelt. Idee: Einzelne
        Bäume überanpassen, aber ihre Fehler sind unkorreliert und mitteln sich weg.
  \item \textbf{Gradient Boosting}: Bäume werden nacheinander gebaut, jeder neue korrigiert die
        Fehler der bisherigen Summe. XGBoost, LightGBM, CatBoost sind optimierte Implementierungen.
  \item \textbf{SVM}: Support Vector Machine, sucht die trennende Hyperebene mit maximalem Abstand
        zu den nächsten Punkten.
  \item Ehrlicher Rechenhinweis aktiv nennen: Der RBF-Kernel-SVM skaliert kubisch mit der
        Zeilenzahl. Bei 1,55 Mio. Trainingszeilen nicht machbar, deshalb auf 8.000 Zeilen gesampelt.
\end{sag}

\zeigen{App $\rightarrow$ \textit{Model Comparison} $\rightarrow$ Abschnitt
\textbf{\enquote{Pareto front: macro-F1 vs. Recall(KSI)}}. Auf den Gate-Bereich deuten und zeigen,
welche Kandidaten ihn erfüllen. Tabelle \enquote{All 10 candidates} nur kurz streifen. 45 Sekunden.}

\rubrik{Begriffe}
\begin{begriff}
  \bg{Bagging vs. Boosting}{Bagging (Random Forest) trainiert viele Modelle parallel auf
      Zufallsstichproben und mittelt. Boosting trainiert sequenziell, jedes Modell korrigiert
      die Fehler der Vorgänger.}
  \bg{Blattweises Wachstum (LightGBM)}{Der Baum wächst dort weiter, wo der Gewinn am größten
      ist, statt Ebene für Ebene. Schneller, aber überanpassungsanfälliger.}
  \bg{Kernel-Trick, RBF}{Erlaubt nichtlineare Trennung, indem implizit in einen
      höherdimensionalen Raum abgebildet wird, ohne diesen explizit zu berechnen. RBF (Radial
      Basis Function) ist der gebräuchlichste nichtlineare Kernel.}
\end{begriff}

% --------------------------------------------------------------------------
\karte{23}{A\textsuperscript{3}}{Champion-Auswahl mit Gate-Logik}

\kernsatz{Random Forest gewinnt als Kandidat mit der höchsten Macro-F1 unter denen, die das
Recall-Gate erfüllen. Das Gate ist Vorbedingung, nicht Zielfunktion.}

\rubrik{Sagen}
\begin{sag}
  \item Die Regel präzise: Nimm alle Kandidaten mit Recall(KSI) von mindestens 0,50. Unter
        diesen nimm den mit der höchsten Macro-F1.
  \item \textbf{Den scheinbaren Widerspruch selbst ansprechen}: Random Forest hat mit 0,540 den
        \textit{niedrigsten} KSI-Recall aller Baum- und SVM-Kandidaten. Trotzdem gewinnt er,
        weil er das Gate erfüllt und unter allen, die es erfüllen, die höchste Macro-F1 hat.
  \item \textbf{Optuna}: Framework für Hyperparameter-Optimierung, das nicht stumpf ein Gitter
        durchprobiert, sondern modellbasiert sucht. Der Tree-structured Parzen Estimator baut aus
        den bisherigen Versuchen ein Wahrscheinlichkeitsmodell darüber, welche Parameterbereiche
        gute Ergebnisse liefern, und probiert dort weiter.
  \item 20 Trials ergaben 180 Bäume, Tiefe 23, mindestens 8 Beobachtungen pro Blatt. Das hob die
        Macro-F1 von 0,6011 auf 0,6083.
\end{sag}

\merksatz{\enquote{Das Gate ist eine Vorbedingung, keine Zielfunktion. Es sortiert aus, es
optimiert nicht.}}

\rubrik{Begriffe}
\begin{begriff}
  \bg{Hyperparameter}{Einstellungen, die vor dem Training festgelegt werden, im Gegensatz zu
      Parametern, die das Modell aus den Daten lernt.}
  \bg{max\_depth}{Maximale Tiefe eines Baums. Begrenzt Komplexität, wirkt gegen Überanpassung.}
  \bg{min\_samples\_leaf}{Mindestanzahl Beobachtungen in einem Blatt. Verhindert, dass Blätter
      auf einzelne Ausreißer gefittet werden.}
\end{begriff}

% --------------------------------------------------------------------------
\karte{24}{A\textsuperscript{3}}{Schwellenwert und Negativbefund \textnormal{\small(kürzbar)}}

\kernsatz{Der Schwellenwert wurde systematisch über 81 Kandidaten optimiert, und ein zweiter
Tuning-Versuch wurde nach einer klaren Regel verworfen statt schöngerechnet.}

\rubrik{Sagen}
\begin{sag}
  \item Punkt, den viele nicht auf dem Schirm haben: Ein Klassifikator gibt eine
        Wahrscheinlichkeit aus. Erst der Schwellenwert macht daraus eine Entscheidung.
        Der Default 0,5 ist eine Konvention, keine Optimierung.
  \item 81 Werte geprüft, der beste unter der Nebenbedingung Recall $\geq$ 0,50: \textbf{0,49860}.
  \item Interessantes Detail: Der optimale Wert liegt sehr nah an 0,5. Kein Zufall, sondern Folge
        der Klassengewichtung, die die Wahrscheinlichkeiten schon vorher grob ausbalanciert hat.
  \item \textbf{Promotionsregel}: Der neue Versuch wird nur übernommen, wenn er auf
        \textit{beiden} Achsen nicht schlechter ist. Der Multi-Objective-Versuch hatte besseren
        Recall (0,5187 gegen 0,5053), aber schlechtere Macro-F1 (0,6057 gegen 0,6083).
        Also verworfen.
\end{sag}

\merksatz{\enquote{Ich hätte auch die Regel nachträglich ändern und den besseren Recall
berichten können. Die Regel stand aber vorher fest, und der Versuch ist als gescheiterte Evidenz
dokumentiert statt gelöscht.}}

\rubrik{Begriffe}
\begin{begriff}
  \bg{Multi-Objective-Optimierung}{Optimierung mehrerer Zielgrößen gleichzeitig. Ergebnis ist
      keine einzelne beste Lösung, sondern eine Pareto-Menge.}
  \bg{Pareto-Front}{Die Menge der Lösungen, bei denen man eine Zielgröße nur verbessern kann,
      indem man eine andere verschlechtert.}
\end{begriff}

% --------------------------------------------------------------------------
\karte{25}{C}{Das Testergebnis 2024}

\kernsatz{Macro-F1 0,604 und Recall(KSI) 0,515 auf einem nie zuvor benutzten Jahr, beide Gates
erfüllt, minimaler Abstand zur Validierung.}

\rubrik{Sagen}
\begin{sag}
  \item Zürst betonen, wie der Test verwendet wurde: genau einmal geöffnet, für genau ein
        Modell, mit dem vorher festgelegten Schwellenwert.
  \item Wichtigster Vergleich, Validierung gegen Test: Macro-F1 0,6083 gegen 0,6039. Beim Recall
        sogar leicht besser, 0,5053 gegen 0,5151.
  \item Ausdrücklich sagen, warum das gut ist: Wäre das Modell auf der Validierung deutlich
        besser als auf dem Test, hätte man den Validierungssatz überangepasst. Hier ist der
        Abstand minimal.
  \item Konfusionsmatrix in Worte fassen: Von rund 44.200 echten KSI-Unfällen wurden 22.767
        erkannt und 21.431 übersehen. Von rund 224.300 Leichtunfällen wurden 172.434 korrekt
        erkannt und 51.887 fälschlich als KSI markiert.
\end{sag}

\merksatz{\enquote{Rund 21.400 KSI-Unfälle werden übersehen. Das Modell ist ein
Priorisierungswerkzeug, kein Sicherheitssystem. Im echten Einsatz müsste die Kostenfunktion des
Anwenders den Schwellenwert bestimmen, nicht Macro-F1. Wenn ein übersehener schwerer Unfall
zehnmal teurer ist als ein Fehlalarm, gehört der Schwellenwert deutlich nach unten.}}

\rubrik{Begriffe}
\begin{begriff}
  \bg{False Negative}{Ein KSI-Unfall, der als leicht eingestuft wurde. Hier der teure Fehler.}
  \bg{False Positive}{Ein Leichtunfall, der als KSI eingestuft wurde. Hier der billigere Fehler,
      kostet nur Aufmerksamkeit.}
\end{begriff}

% --------------------------------------------------------------------------
\karte{26}{C}{Die Entscheidungsmatrix}

\kernsatz{Über vier gewichtete Kriterien führt XGBoost, nicht das ausgelieferte Modell.
Dieses unbequeme Ergebnis wird berichtet statt versteckt.}

\rubrik{Sagen}
\begin{sag}
  \item Methodik zürst: Sechs Kriterien waren nominal vorgesehen, jedes mit einem Gewicht. Jedes
        wird min-max-normalisiert, also auf 0 bis 1 skaliert. Kostenkriterien wie Latenz werden
        invertiert, weil dort niedriger besser ist.
  \item \textbf{Methodische Sorgfalt zeigen}: Zwei Kriterien, Interpretierbarkeit und
        Trainingskosten, wurden nie gemessen. Statt sie subjektiv zu schätzen und damit das
        Ergebnis zu färben, wurden sie automatisch ausgeschlossen und die restlichen Gewichte
        renormalisiert.
  \item Ergebnis offensiv nennen: XGBoost führt mit 0,561, Random Forest ist Letzter mit 0,500.
        XGBoost hat mehr Recall und ist dreieinhalbmal schneller.
  \item Ankündigen, dass die nächste Folie erklärt, warum trotzdem nicht gewechselt wird.
        Das erzeugt Spannung und verhindert eine vorschnelle Frage.
\end{sag}

\rubrik{Begriffe}
\begin{begriff}
  \bg{Min-Max-Normalisierung}{Skalierung auf 0 bis 1 über (Wert minus Minimum) geteilt durch
      (Maximum minus Minimum).}
  \bg{Kostenkriterium}{Ein Kriterium, bei dem kleiner besser ist, hier die Latenz. Wird invertiert,
      damit größer immer besser bedeutet.}
  \bg{Latenz ms/1k}{Millisekunden für 1.000 Vorhersagen.}
\end{begriff}

% --------------------------------------------------------------------------
\karte{27}{C}{Warum der Champion Champion bleibt \textnormal{\small(Schlüsselfolie)}}

\kernsatz{Den Champion jetzt zu wechseln würde den Testsatz in den Auswahlprozess ziehen und
damit genau die Unabhängigkeit zerstören, die ihn wertvoll macht.}

\rubrik{Vier Schritte, 90 Sekunden, wichtigste Folie der Präsentation}
\begin{sag}
  \item \textbf{1. Ausgangslage}: Die Modellauswahl fiel in Phase A\textsuperscript{3} auf Basis
        der Validierungsdaten 2023. Zu diesem Zeitpunkt war Test 2024 unberührt, reserviert für
        genau einen Zweck: eine bereits eingefrorene Entscheidung zu bestätigen.
  \item \textbf{2. Was ein Wechsel bedeuten würde}: Wollte ich XGBoost ausliefern, müsste ich
        für ihn eine Testzahl berichten. Also müsste ich Test 2024 auch für XGBoost auswerten.
  \item \textbf{3. Warum das den Test entwertet}: Dann ist der Testsatz Teil der Auswahl geworden.
        Ich habe zwei Modelle auf dem Test gemessen und das bessere genommen. Die berichtete Zahl
        ist das Maximum aus zwei Versuchen, keine unverzerrte Schätzung der
        Generalisierungsleistung mehr. Rückwirkend wäre auch die Random-Forest-Zahl nicht mehr
        unabhängig.
  \item \textbf{4. Der Name des Problems}: Data Leakage auf der Ebene der Modellauswahl. Subtiler
        als Leakage auf Merkmalsebene, weil kein einziger Datenpunkt ins Training wandert.
        Trotzdem sickert Information aus dem Test in die Entscheidung. In der Praxis die
        häufigere Form.
  \item Konsequenz: Random Forest bleibt für diesen Zyklus. XGBoost ist dokumentiert als Kandidat
        für einen zukünftigen, vorab registrierten Vergleich.
\end{sag}

\merksatz{\enquote{Die Validierung entscheidet. Der Test bestätigt nur.}}

\zeigen{App $\rightarrow$ \textit{Model Comparison} $\rightarrow$ ganz unten den Aufklapper
\textbf{\enquote{Warum bleibt Random Forest der Champion? (German explanation)}} öffnen. Nur
zeigen, dass die Überlegung im Produkt steht, nicht vorlesen. 30 Sekunden.}

\rubrik{Wenn gefragt}
\begin{begriff}
  \frage{Warum nicht einfach einen neuen Testsatz nehmen?}{Es gibt keinen, 2024 ist das letzte
    verfügbare Jahr. Sobald 2025 veröffentlicht ist, wäre genau das der saubere Weg: vorab
    festlegen, dass XGBoost gegen Random Forest auf 2025 verglichen wird, und dann einmal messen.}
  \frage{Ist das nicht übertrieben streng?}{Es ist der Standard, der reproduzierbare Forschung
    von Zahlenpolitur trennt. Der ganze Wert eines Holdout-Sets liegt darin, dass es genau einmal
    benutzt wird.}
\end{begriff}

% --------------------------------------------------------------------------
\karte{28}{C}{Permutation Importance}

\kernsatz{Unfallart und Unfalltyp dominieren, gefolgt von der Beteiligung ungeschützter
Verkehrsteilnehmer. Modellagnostisch gemessen, damit alle vier Finalisten vergleichbar sind.}

\rubrik{Sagen}
\begin{sag}
  \item Verfahren zürst, es ist intuitiv: Nimm eine Merkmalsspalte und mische ihre Werte über
        alle Zeilen zufällig durch. Der Zusammenhang zum Ziel ist zerstört, die Verteilung der
        Spalte bleibt unverändert. Dann miss, wie stark die Modellleistung einbricht. Großer
        Einbruch bedeutet: Das Modell war auf dieses Merkmal angewiesen.
  \item Vorteil gegenüber der eingebauten Feature Importance von Bäumen: modellagnostisch. Man
        kann Random Forest, XGBoost, LightGBM und CatBoost auf derselben Stichprobe und mit
        derselben Metrik vergleichen.
  \item Rangliste plausibilisieren: Unfallart und Unfalltyp vorn. Ein Frontalzusammenstoß ist
        strukturell gefährlicher als ein Auffahrunfall im Stau. Dann Motorrad- und
        Fahrradbeteiligung, also ungeschützte Verkehrsteilnehmer. Das bestätigt H2 in der
        schwachen Form.
  \item \textbf{Warum nicht SHAP?} SHAP erklärt, wie ein \textit{einzelnes} Modell seine
        transformierten Merkmale nutzt. Für einen Vergleich \textit{zwischen} Modellen das falsche
        Werkzeug, weil jedes Modell einen anderen transformierten Merkmalsraum hat. Permutation
        Importance arbeitet auf den Rohspalten und ist deshalb vergleichbar.
  \item Ehrlicher Zusatz, falls Zeit ist: Die Rangordnungen der vier Modelle stimmen nur mässig
        überein. Mittlere Spearman-Korrelation zum Champion 0,37, Top-10-Überschneidung etwa
        40 Prozent. Die Modelle sind sich einig, \textit{dass} Unfallart wichtig ist, aber uneinig
        über die zweite Reihe. Die Rangliste also nicht überinterpretieren.
\end{sag}

\rubrik{Begriffe}
\begin{begriff}
  \bg{SHAP}{SHapley Additive exPlanations. Verteilt die Abweichung einer Vorhersage vom Mittelwert
      additiv auf die Merkmale, basierend auf dem Shapley-Wert aus der kooperativen Spieltheorie.}
  \bg{Spearman-Korrelation}{Korrelation der Ränge statt der Werte. Passend beim Vergleich von
      Rangordnungen.}
  \bg{Jaccard-Index}{Größe der Schnittmenge geteilt durch Größe der Vereinigungsmenge.
      Hier: Wie stark überlappen zwei Top-10-Listen?}
\end{begriff}

% --------------------------------------------------------------------------
\karte{29}{C}{Robustheit und Modell-Uneinigkeit}

\kernsatz{Random Forest ist mit 3,77 Prozent gekippten Vorhersagen das stabilste Modell bei
fehlenden Merkmalen, und die Boosting-Modelle sind sich untereinander deutlich einiger als mit ihm.}

\rubrik{Sagen}
\begin{sag}
  \item Robustheitssonde erklären: Fünf Merkmale, die im echten Betrieb realistisch ausfallen
        können, werden einzeln komplett auf NaN gesetzt. Also: Was passiert, wenn die
        OSM-Anreicherung ausfällt oder die Wetterstation keine Daten liefert? Gemessen wird,
        wie viele Vorhersagen dadurch ihre Klasse wechseln.
  \item Ergebnis: Random Forest 3,77 Prozent, die Boosting-Modelle zwischen 7 und 9 Prozent.
  \item Eindrücklicher Vergleichswert: SVM über 52 Prozent gekippte Vorhersagen bei einem
        einzigen fehlenden Merkmal. Starkes praktisches Argument für Baumensembles im Betrieb.
  \item Paarweise Uneinigkeit: Random Forest gegen die Boosting-Modelle rund 13 Prozent,
        die Boosting-Modelle untereinander nur 5 Prozent.
  \item Interpretation: Die drei Boosting-Modelle sind faktisch Varianten derselben Idee. Random
        Forest ist der methodische Ausreißer, er optimiert auf ausgewogene Klassenleistung und
        drängt nicht so stark auf KSI-Recall.
  \item Schlussfolgerung anbieten: Für ein Ensemble wäre Random Forest plus ein Boosting-Modell
        interessanter als zwei Boosting-Modelle, weil unkorrelierte Fehler sich besser ausgleichen.
\end{sag}

\rubrik{Begriffe}
\begin{begriff}
  \bg{NaN}{Not a Number, die Standardkodierung für fehlende Werte.}
  \bg{Paarweise Uneinigkeit}{Anteil der Fälle mit unterschiedlicher harter Klassenzuweisung
      zwischen zwei Modellen.}
\end{begriff}

% --------------------------------------------------------------------------
\karte{30}{C}{Limitationen}

\kernsatz{Die Leistungsgrenze liegt in den Daten, nicht im Modell, und alle wesentlichen
Einschränkungen sind dokumentiert.}

\rubrik{Sagen}
\begin{sag}
  \item Zügig, aber vollständig. Das ist eine Stärke, keine Schwäche. Offen benannte
        Limitationen werden positiv bewertet.
  \item Wichtigster Punkt sind die fehlenden Merkmale. Zurückverweisen auf Folie 16 und 20:
        Das ist die Ursache der Leistungsdecke, und sie ist nicht durch bessere Modellierung
        behebbar.
  \item Konzeptionell interessanter Schluss: Der Schwellenwert ist eine Wertentscheidung, keine
        technische. Er kodiert implizit, wie teuer ein übersehener schwerer Unfall im Vergleich
        zu einem Fehlalarm ist. Diese Abwägung gehört dem Anwender, nicht dem Entwickler.
\end{sag}

\rubrik{Wenn gefragt}
\begin{begriff}
  \frage{Was würde die Leistung am meisten verbessern?}{Eindeutig zusätzliche Merkmale, keine
    besseren Algorithmen. Am wertvollsten wären Aufprallgeschwindigkeit und Insassenalter. Beide
    sind in Unfallforschungsdatenbanken vorhanden, aber nicht im offenen Datensatz, weil sie
    personenbezogen oder erhebungsaufwendig sind.}
\end{begriff}

% --------------------------------------------------------------------------
\karte{31}{C}{Der Inference Contract \textnormal{\small(kürzbar)}}

\kernsatz{Eine maschinenlesbare Vertragsdatei fixiert Schema, Wertebereiche, Schwellenwert und
Artefakt-Prüfsumme, damit die App ohne Notebook-Ausführung korrekt arbeitet.}

\rubrik{Sagen}
\begin{sag}
  \item Problem: Wer eine Anwendung auf ein Modell aufsetzt, muss wissen, welche Spalten in
        welchem Typ und Wertebereich erwartet werden, welcher Schwellenwert gilt und welche
        Modelldatei die richtige ist. Normalerweise steht das verstreut in Notebooks und wird
        beim Nachbau falsch rekonstruiert.
  \item Lösung: eine JSON-Datei, die all das enthält, generiert aus den echten Trainingsdaten,
        nicht aus Annahmen. 30 Pflichtspalten mit Typ, Herkunft und Wertebereich.
  \item \textbf{Stärkster Einzelpunkt, SHA-256}: Beim Laden wird die Prüfsumme jeder Modelldatei
        neu berechnet und gegen den Eintrag verglichen. Passt sie nicht, bricht die Analyse ab.
        Damit ist ausgeschlossen, dass die Metadaten ein Modell beschreiben, das inzwischen ein
        anderes ist.
\end{sag}

\rubrik{Begriffe}
\begin{begriff}
  \bg{SHA-256}{Kryptografische Hashfunktion. Erzeugt aus beliebigem Inhalt einen Fingerabdruck
      fester Länge. Schon eine minimale Änderung am Inhalt ändert den Hash vollständig.}
  \bg{Schema}{Die formale Beschreibung, welche Spalten mit welchen Typen erwartet werden.}
\end{begriff}

% --------------------------------------------------------------------------
\karte{32}{K}{Die Streamlit-Anwendung}

\kernsatz{Vier Seiten, die jeweils eine konkrete Nutzerfrage beantworten, lauffähig allein aus
eingecheckten Artefakten.}

\rubrik{Sagen}
\begin{sag}
  \item Die vier Seiten kurz benennen, jeweils mit der Frage, die sie beantwortet. Nicht alle
        Details vorlesen.
  \item Wichtigster technischer Punkt: Die App braucht keine Notebook-Ausführung. Alle Artefakte,
        also Modell, Contract, Kennzahlen und Kartendaten, liegen im Repository. Ein Prüfer kann
        klonen und starten.
\end{sag}

\zeigen{App $\rightarrow$ \textit{Risk Predictor}. In dieser Reihenfolge:
\begin{sag}[itemsep=0.9pt]
  \item Abschnitt \textbf{\enquote{Try an example}}, ein Szenario anklicken. Das füllt alle Felder.
  \item Kurz durch \textbf{\enquote{Pick a location}} (Kartenklick) und
        \textbf{\enquote{Accident characteristics}} scrollen, nicht erklären.
  \item \textbf{\enquote{Predict KSI risk}} drücken, Wahrscheinlichkeit und Schwellenwert zeigen.
  \item Optional: Seite \textit{Why This Prediction}, Balkendiagramm der Merkmalswichtigkeiten.
\end{sag}
Rückfalloption: Startet die App nicht, bleibt diese Folie stehen und ich rede einfach weiter.
60 Sekunden.}

\rubrik{Begriffe}
\begin{begriff}
  \bg{Streamlit}{Python-Framework, das aus einem Skript eine Web-App macht. Bei jeder Interaktion
      läuft das Skript komplett neu, Caching verhindert die Wiederholung teurer Berechnungen.}
  \bg{Artefakt}{Eine gespeicherte Datei als Ergebnis eines Analyseschritts, hier Modelle,
      Kennzahlen-CSVs und die Contract-JSON.}
\end{begriff}

% --------------------------------------------------------------------------
\karte{33}{K}{Die Risikokarte und die Glättung}

\kernsatz{Statt absoluter Anteile zeigt die Karte relatives Risiko mit Bayes'scher Glättung,
weil sonst 97,5 Prozent der Fläche einfarbig wären.}

\rubrik{Sagen}
\begin{sag}
  \item Problem zürst, es ist anschaulich: KSI macht bundesweit 18,91 Prozent aus. Färbt man
        eine Zelle erst rot, wenn die Mehrheit der Unfälle KSI ist, bräuchte sie das 2,64-Fache
        des Durchschnitts. Das schaffen 123 von 4.857 Zellen. Die Karte wäre praktisch einfarbig.
  \item Lösung: Nicht \enquote{ist die Mehrheit schwer?}, sondern \enquote{ist diese Zelle
        schlechter als der Durchschnitt, und um wie viel?}.
  \item Zweites Problem, das Glättung nötig macht: Eine Zelle mit drei Unfällen, davon zwei KSI,
        hätte eine Rate von 67 Prozent, also das 3,5-Fache des Durchschnitts. Das ist reines
        Rauschen bei drei Beobachtungen.
  \item \textbf{Shrinkage}: Man tut so, als hätte man zusätzlich k Pseudo-Unfälle beobachtet,
        die exakt dem Bundesdurchschnitt entsprechen. Bei vielen echten Beobachtungen fallen sie
        kaum ins Gewicht, bei wenigen ziehen sie die Schätzung stark zum Durchschnitt. Im Kern
        eine Bayes'sche Schätzung mit dem Bundesdurchschnitt als Prior.
  \item \textbf{Warum k gleich 20}, empirisch bestimmt, nicht geraten: Bei k gleich 10 rutschen
        noch 6-Unfall-Zellen ins höchste Band, bei k gleich 50 bleiben nur 52 Zellen übrig und
        die Karte wird flach. Bei 20 sind alle fünf Bänder sinnvoll besetzt.
  \item Schöner UX-Punkt: Zellen mit wenig Evidenz werden blasser gezeichnet. Der Nutzer sieht
        nicht nur den Schätzwert, sondern auch, wie sicher er ist.
\end{sag}

\zeigen{App $\rightarrow$ \textit{Overview} $\rightarrow$ Abschnitt \textbf{\enquote{Where
accidents are disproportionately severe}}. Ein oberes Risikoband aus- und wieder einschalten,
damit man sieht, wie stark es ausdünnt. Auf eine blasse Zelle deuten und die Konfidenz-Codierung
erklären. 45 Sekunden.}

\rubrik{Wenn gefragt}
\begin{begriff}
  \frage{Ist das nicht Manipulation der Daten?}{Nein, eine bewusste dokumentierte Verzerrung
    zugunsten der Robustheit. Ohne Glättung wären die auffälligsten Zellen systematisch die mit
    den wenigsten Beobachtungen, weil kleine Stichproben die extremsten Raten produzieren. Bei
    Karten seltener Ereignisse ist das in der Epidemiologie Standard.}
\end{begriff}

% --------------------------------------------------------------------------
\karte{34}{K}{Engineering und Qualitätssicherung}

\kernsatz{Ein Fehler, den die Standard-Tests nicht sehen konnten, führte zu einem zusätzlichen
Browser-Test als Pflicht-Gate.}

\rubrik{Sagen, die Geschichte kurz erzählen, sie bleibt hängen}
\begin{sag}
  \item Die Kartenbibliothek verändert die Objekte, die man ihr übergibt, während sie rendert.
        Speichert man diese Objekte zwischen, um Zeit zu sparen, referenziert ein späterer Aufruf
        eine JavaScript-Variable, die es nicht mehr gibt. Die Karte bleibt weiß.
  \item \textbf{Der tückische Teil}: Streamlits eigenes Test-Framework führt keinen Browser aus.
        Es hat null Fehler gemeldet, während die Karte im Browser kaputt war. Der Test war grün,
        das Produkt war defekt.
  \item Konsequenz: ein Headless-Browser-Test, der eine echte Browser-Instanz startet und auf
        JavaScript-Fehler prüft. Plus ein Test, der die Invariante absichert, dass diese Objekte
        nie zwischengespeichert werden.
  \item Restliche Qualitätssicherung kurz: über 440 Tests, Linting, Pre-Commit, CI bei jedem
        Push, Notebooks im Audit-Modus.
\end{sag}

\merksatz{\enquote{Ein grüner Test ist nur so viel wert wie das, was er tatsächlich ausführt.}}

\rubrik{Begriffe}
\begin{begriff}
  \bg{Headless Browser}{Ein echter Browser ohne sichtbare Oberfläche, für automatisierte Tests.}
  \bg{CI}{Continuous Integration. Automatische Ausführung von Tests und Prüfungen bei jeder
      Änderung im Repository.}
  \bg{Linting}{Statische Codeprüfung auf Stil- und Fehlerpotenzial. \textbf{Pre-Commit-Hook}:
      Prüfung, die vor jedem Commit automatisch läuft.}
\end{begriff}

% --------------------------------------------------------------------------
\karte{35}{Schluss}{Fazit}

\kernsatz{Das Modell funktioniert, aber der eigentliche Beitrag liegt in der methodischen
Disziplin an drei Stellen, an denen sie etwas gekostet hat.}

\rubrik{Sagen}
\begin{sag}
  \item Ergebnis nennen: 0,604 und 0,515, beide Gates erfüllt.
  \item Dann die drei Punkte, jeweils mit der Betonung, dass sie etwas gekostet haben:
  \begin{sag}
    \item \textbf{Der Negativbefund wurde nicht versteckt.} Es wäre einfacher gewesen, die
          Drei-Klassen-Ergebnisse wegzulassen und nur das binäre Modell zu zeigen.
    \item \textbf{Das Reframing war vorab definiert.} Das hat Arbeit gemacht in Phase Q, als noch
          niemand wusste, dass man es brauchen würde.
    \item \textbf{Die Test-Trennung wurde eingehalten, obwohl sie wehtat.} Die Entscheidungsmatrix
          bevorzugt XGBoost. Es wäre attraktiver gewesen, das bessere Modell zu präsentieren.
  \end{sag}
  \item Ausblick kurz, vier Stichpunkte, nicht ausführen.
\end{sag}

\merksatz{\enquote{Was die Grenze setzt, sind die Daten, nicht das Modell. Kein Algorithmus kann
Information erzeugen, die nicht in den Daten steckt.}}

% --------------------------------------------------------------------------
\karte{36}{Schluss}{Quellen}

\rubrik{Sagen}
\begin{sag}
  \item Nicht vorlesen. Ein Satz genügt: Datensätze oben, Methodenquellen unten, alle im
        APA-7-Stil.
\end{sag}

\rubrik{Wenn nach Fachliteratur zur Unfallforschung gefragt wird}
\merksatz{\enquote{Ich zitiere hier bewusst nur Quellen, die ich vollständig belegen kann, also
die Datensätze und die Methodenpapiere zu den eingesetzten Verfahren. Für den Vergleich mit
publizierten Ergebnissen aus der Unfallforschung gilt ohnehin eine Einschränkung, die im
Repository dokumentiert ist: Solche Studien arbeiten meist mit Personen- und Fahrzeugmerkmalen,
die dieser Datensatz nicht hat, oder mit einem anderen Zielformat. Ein direkter Zahlenvergleich
wäre deshalb irreführend.}}

\rubrik{Warum das eine starke Antwort ist}
\begin{sag}
  \item Sie benennt die Grenze der Vergleichbarkeit selbst, statt sie zu verschweigen.
\end{sag}

% --------------------------------------------------------------------------
\karte{37}{Schluss}{Vielen Dank}

\rubrik{Sagen}
\begin{sag}
  \item Kurz halten.
  \item Auf die drei Artefakte verweisen: Repository, Live-Report, Streamlit-App.
  \item Fragen einladen.
  \item Die App im Hintergrund offen lassen, um bei Fragen direkt hineinspringen zu können.
\end{sag}

\zeigen{Bei Fragen zur Umsetzung: App $\rightarrow$ \textit{Why This Prediction} zeigt für die
zuletzt berechnete Vorhersage die global wichtigsten Merkmale mit den eigenen Eingabewerten.
Gute Antwort auf \enquote{Woran sieht man, was das Modell benutzt hat?}.}

% ============================================================== NOTFALLKARTE
\metakarte{Notfall}{Die härteste Rückfrage in je einem Satz}

\begin{begriff}
  \frage{Warum nicht Accuracy?}{81 Prozent Accuracy bekommt man geschenkt, wenn man immer die
    häufigste Klasse sagt. Macro-F1 gewichtet jede Klasse gleich.}
  \frage{Warum ist das Modell so schwach?}{Weil die physikalischen Determinanten fehlen:
    Aufprallgeschwindigkeit, Insassenalter, Gurtnutzung, Fahrzeugmasse. Die Grenze liegt in den
    Daten, nicht im Algorithmus.}
  \frage{Haben Sie das Ziel geändert, weil es nicht funktioniert hat?}{Ja, und das ist
    dokumentiert. Entscheidend ist: Der Ausweichpfad stand in Phase Q, bevor das erste Modell
    trainiert wurde, und die ursprüngliche Zielverfehlung wird vollständig berichtet.}
  \frage{Warum liefern Sie nicht das bessere Modell aus?}{Weil ich dafür den Testsatz ein zweites
    Mal benutzen müsste. Damit wäre er Teil der Auswahl und keine unabhängige Schätzung mehr.
    Die Validierung entscheidet, der Test bestätigt nur.}
  \frage{Ist Cram\'ers V 0,13 nicht zu wenig für ein Modell?}{Für die Drei-Klassen-Aufgabe ja,
    das ist der berichtete Negativbefund. Für die binäre KSI-Aufgabe reicht es, weil die
    positive Klasse dort 19 statt 1 Prozent ausmacht.}
  \frage{Kann man damit einzelne Personen bewerten?}{Nein, und das war eine vorab festgelegte
    Ausschlussentscheidung. Der Datensatz enthält keine Personendaten, und die Ergebnisse
    informieren Infrastrukturentscheidungen.}
  \frage{Sind die Merkmalswichtigkeiten Ursachen?}{Nein, Assoziationen. Dieser Hinweis steht auch
    in der App.}
  \frage{Warum kein neuronales Netz?}{Interpretierbarkeit war eine harte Vorgabe. Ein
    Verkehrsplaner muss nachvollziehen können, warum ein Ort als Risiko markiert wurde.}
\end{begriff}

\end{document}
